\subsection{Orden Topológico (Kahn) y DP sobre DAG con Reconstrucción}

\graybold{Preguntas guía:}

\begin{itemize}
\item ¿El grafo es dirigido y SIN ciclos (DAG), o debo detectar si tiene ciclos?
\item ¿Me piden el camino más largo (o más corto, o la cantidad de caminos) entre dos nodos de un DAG?
\item ¿El valor de cada nodo depende solo de los valores de sus predecesores, de modo que basta procesarlos en un orden donde todo predecesor va antes?
\item ¿Además del valor óptimo debo imprimir el camino, guardando de qué nodo vino cada mejora?
\end{itemize}

\graybold{Utilidad:} El orden topológico ordena los nodos de un DAG de forma que toda arista va de un nodo anterior a uno posterior. Recorriendo los nodos en ese orden, cuando se llega a un nodo todos sus predecesores ya tienen su valor definitivo, así que cualquier DP que dependa solo de ellos se calcula en una pasada: camino más largo o más corto, cantidad de caminos, tareas con dependencias (tiempo mínimo de finalización), planificación de cursos. Kahn además detecta ciclos: si no logra ordenar todos los nodos, el grafo tiene uno.

\graybold{Complejidad:} \bigO{n + m}: cada nodo entra una vez a la cola y cada arista se revisa una vez.

\graybold{Fundamento teórico:} El algoritmo de Kahn mantiene el grado de entrada de cada nodo y una cola con los que tienen grado 0, es decir, sin predecesores pendientes. Al procesar un nodo, se "borran" sus aristas salientes restando 1 al grado de cada sucesor, y los que llegan a 0 entran a la cola. Un nodo entra a la cola exactamente cuando todos sus predecesores fueron procesados, así que el orden de la cola es topológico. Si el grafo tiene un ciclo, ningún nodo del ciclo llega nunca a grado 0 (cada uno espera al anterior), y la cola termina con menos de $n$ nodos: esa es la detección de ciclos. Sobre ese orden, el DP del camino más largo es $\text{dp}[v] = \max(\text{dp}[u] + 1)$ sobre las aristas $u \to v$, y es correcto porque al procesar $u$ su valor ya es definitivo. En este problema los caminos deben EMPEZAR en la ciudad 1, así que $\text{dp}[v] = 0$ significa "inalcanzable desde 1" y solo se propagan valores desde nodos con $\text{dp} > 0$; aun así Kahn tiene que recorrer todo el grafo, porque los nodos inalcanzables también descuentan grados de entrada, y sin ellos otros nodos nunca llegarían a 0. La ciudad 1 puede tener aristas entrantes desde nodos inalcanzables y por eso no estar al principio de la cola, pero ningún nodo procesado antes que ella es alcanzable desde ella en un DAG, así que su valor inicial 1 nunca se pierde. Para reconstruir, cada vez que $\text{dp}[v]$ mejora se guarda el predecesor responsable, y desde el destino se sigue esa cadena hasta el origen.

\graybold{Ejercicio aplicado}

Dado un DAG de $n$ ciudades y $m$ vuelos dirigidos, hallar la ruta de la ciudad 1 a la ciudad $n$ que visita la mayor cantidad de ciudades e imprimirla, o informar que no existe.

\graybold{I/O:} $n \le 10^5$, $m \le 2\cdot10^5$ con dos enteros por arista $\Rightarrow$ lectura total + tokenización (patrón 2), separando extremos con slices de paso 2. Como se acepta cualquier ruta óptima, la verificación comprueba que la ruta empiece en 1, termine en $n$, use solo vuelos existentes y tenga la longitud máxima, calculada enumerando TODAS las rutas: se hizo en 600 DAGs aleatorios, incluyendo la ciudad 1 con aristas entrantes, $n$ inalcanzable y varias rutas óptimas, sin fallos. Con $n = 10^5$ y $m = 2\cdot10^5$ tarda entre \aprox{}0.08s y \aprox{}0.19s: una ruta de \ten{5} ciudades, DAGs aleatorios con y sin camino a $n$ (en uno de ellos la respuesta IMPOSSIBLE se confirmó con un recorrido independiente), y $n$ inalcanzable; el límite es de 1 segundo. El toposort genérico de los extras se verificó contra un detector de ciclos por DFS en 2000 grafos con y sin ciclos (incluyendo lazos), comprobando además que cada orden devuelto respete todas las aristas, y tarda \aprox{}0.04s con \ten{5} nodos.

\graybold{Código solución}

\begin{lstlisting}[language=Python]
import sys

def solve():
    data = sys.stdin.buffer.read().split()
    n = int(data[0]); m = int(data[1])
    adj = [[] for _ in range(n + 1)]
    grado = [0]*(n + 1)                   # grado de entrada
    for a, b in zip(map(int, data[2:2 + 2*m:2]), map(int, data[3:3 + 2*m:2])):
        adj[a].append(b)
        grado[b] += 1
    dp = [0]*(n + 1)                      # ciudades en la mejor ruta desde 1; 0 = inalcanzable
    dp[1] = 1
    padre = [0]*(n + 1)
    cola = [v for v in range(1, n + 1) if grado[v] == 0]
    for u in cola:                        # la lista crece mientras se recorre: orden topologico
        du = dp[u]
        if du:                            # alcanzable desde 1: propaga el DP
            du1 = du + 1
            for v in adj[u]:
                if du1 > dp[v]:
                    dp[v] = du1
                    padre[v] = u          # de donde vino la mejora, para reconstruir
                grado[v] -= 1
                if not grado[v]:
                    cola.append(v)
        else:                             # inalcanzable: igual descuenta grados,
            for v in adj[u]:              # o sus sucesores nunca entrarian a la cola
                grado[v] -= 1
                if not grado[v]:
                    cola.append(v)
    if not dp[n]:
        sys.stdout.write("IMPOSSIBLE\n"); return
    camino = []
    v = n
    while v:                              # padre[1] = 0 corta la cadena
        camino.append(v)
        v = padre[v]
    camino.reverse()
    sys.stdout.write("%d\n%s\n" % (len(camino), " ".join(map(str, camino))))

solve()
\end{lstlisting}

\graybold{Extras: orden topológico genérico con detección de ciclos} — la parte reutilizable de Kahn, separada del DP. Devuelve el orden, o \texttt{None} si el grafo tiene un ciclo; cualquier DP sobre el DAG se hace luego recorriendo \texttt{orden} (o \texttt{reversed(orden)} si el valor de un nodo depende de sus sucesores).

\begin{lstlisting}[language=Python]
def toposort(n, adj):
    grado = [0]*(n + 1)
    for u in range(1, n + 1):
        for v in adj[u]:
            grado[v] += 1
    orden = [v for v in range(1, n + 1) if grado[v] == 0]
    for u in orden:                       # la lista crece mientras se recorre
        for v in adj[u]:
            grado[v] -= 1
            if not grado[v]:
                orden.append(v)
    # si hay un ciclo, sus nodos nunca llegan a grado 0 y quedan afuera
    return orden if len(orden) == n else None
\end{lstlisting}
